\renewcommand{\Anno}{2000}

\renewcommand{\NumeroProblema}{4}

\renewcommand{\Testo}{La serratura della valigia di Anna è dotata
di una combinazione composta da cinque cilindri affiancati, ognuno dei quali può assumere ciclicamente 10 posizioni, contrassegnate nell'ordine dalle lettere \texttt{A}, \texttt{B}, \texttt{C},
\texttt{D}, \texttt{E}, \texttt{F}, \texttt{M}, \texttt{N},
\texttt{O}, \texttt{P}.
	
Il congegno dovrebbe permettere lo spostamento di ogni singolo cilindro dalla lettera su cui è posizio\-na\-to alla lettera successiva o precedente.
	
Esso però non è perfettamente funzionante, e quando viene spostato di una posizione un qualsiasi cilindro, allora anche uno (e solo uno) di quelli adiacenti si sposta di una posizione nello stesso verso (per esempio supponiamo che la serratura sia posizionata sulla parola \texttt{NONNA}, e che si sposti il quarto cilindro dalla
\texttt{N} alla \texttt{M}: allora anche il terzo cilindro si sposta
dalla \texttt{N} alla \texttt{M}, o il quinto cilindro si sposta dalla
\texttt{A} alla \texttt{P}).  Ovviamente se ad essere spostati sono il
primo od il quinto cilindro, allora a muoversi saranno il secondo od il quarto, rispettivamente.
	
Attualmente la serratura è posizionata sulla parola \texttt{AAAAA}.
\begin{itemize}
	\item Sapendo che la serratura si libera solo sulla parola
	\texttt{BABBO}, riuscirà Anna ad aprire la valigia?
	
	\item E se la serratura si liberasse solo sulla parola
	\texttt{MAMMA}?
\end{itemize}}

\renewcommand{\Risposta}{In entrambi i casi è impossibile aprire
la valigia.}

\renewcommand{\Aiutino}{Intanto conviene matematizzare il problema
sostituendo alle lettere i numeri 1, 2, \ldots, 9, 0.  Ora non resta che trovare degli invarianti rispetto alle mosse possibili.  Dette $a$, $b$, $c$, $d$, $e$ le 5 cifre della combinazione, due invarianti
utili sono $a+b+c+d+e$ modulo \ldots\ e $a-b+c-d+e$ modulo \ldots.}

\renewcommand{\SoluzioneUfficiale}{
}

\renewcommand{\SoluzioniAlternative}{
}

\renewcommand{\Osservazioni}{
}

\renewcommand{\Commenti}{
}
% Aggiunta redazionale: la soluzione preesistente non viene modificata.
\providecommand{\SoluzioneRedazionale}{}
\renewcommand{\SoluzioneRedazionale}{
% SSSUP-BOZZA-AI: SoluzioneRedazionale; autore=ChatGPT; data=2026-11-08; stato=DA-VERIFICARE
% SSSUP-PROVENIENZA: AI: elaborazione nuova dal testo
Codifichiamo le lettere $A,B,C,D,E,F,M,N,O,P$ mediante i residui $0,1,2,3,4,5,6,7,8,9$ modulo $10$. Ogni mossa aggiunge $1$ oppure $-1$ a due posizioni adiacenti. Di conseguenza la quantità
\[
I=x_1-x_2+x_3-x_4+x_5\pmod{10}
\]
è invariante: i contributi delle due posizioni modificate si cancellano. Questo ragionamento comprende anche i passaggi ciclici tra $P$ e $A$.

All'inizio, nella configurazione $AAAAA$, l'invariante è zero. Per $BABBO$ vale invece
\[
1-0+1-1+8\equiv9\pmod{10},
\]
mentre per $MAMMA$ vale
\[
6-0+6-6+0\equiv6\pmod{10}.
\]
Nessuna delle due combinazioni è raggiungibile, qualunque siano le scelte dei cilindri adiacenti che si muovono insieme a quello azionato.
}
